% ============================================================================
% WORKBOOK — Additional: Laws of Exponents
% State-specific (FL, VA)
% Practice-focused reinforcement for the study guide topic
% ============================================================================
\section{Laws of Exponents}
\topicTitle{Laws of Exponents}

\begin{quickReview}{Exponent Rules}
\begin{center}
\begin{tabular}{l l l}
    \textbf{Rule} & \textbf{Formula} & \textbf{Example} \\
    \hline
    Product Rule & $a^m \cdot a^n = a^{m+n}$ & $2^3 \cdot 2^4 = 2^7 = 128$ \\
    Quotient Rule & $\dfrac{a^m}{a^n} = a^{m-n}$ & $\dfrac{5^6}{5^2} = 5^4 = 625$ \\
    Power of a Power & $(a^m)^n = a^{mn}$ & $(3^2)^3 = 3^6 = 729$ \\
    Zero Exponent & $a^0 = 1$ \;($a \neq 0$) & $7^0 = 1$ \\
\end{tabular}
\end{center}

\medskip
\textbf{Key:} Exponent rules only work when the \textbf{bases are the same}. $2^3 \cdot 3^2$ cannot be simplified with these rules.
\end{quickReview}

% ── Warm-Up ──────────────────────────────────────────────────────────────────
\begin{practiceBox}[funGreen]{\faSun~Warm-Up}

\practiceHeader[funGreen]{Evaluate Each Power}

\begin{multicols}{2}
\prob $3^4 =$ \answerBlank[2cm]
\answerExplain{$81$}{$3^4 = 3 \times 3 \times 3 \times 3 = 81$.}

\prob $2^5 =$ \answerBlank[2cm]
\answerExplain{$32$}{$2^5 = 2 \times 2 \times 2 \times 2 \times 2 = 32$.}

\prob $5^3 =$ \answerBlank[2cm]
\answerExplain{$125$}{$5^3 = 5 \times 5 \times 5 = 125$.}

\prob $10^4 =$ \answerBlank[2cm]
\answerExplain{$10{,}000$}{$10^4 = 10 \times 10 \times 10 \times 10 = 10{,}000$.}

\prob $(-2)^3 =$ \answerBlank[2cm]
\answerExplain{$-8$}{$(-2)^3 = (-2)(-2)(-2) = -8$. A negative base raised to an odd power is negative.}

\prob $\left(\frac{1}{3}\right)^2 =$ \answerBlank[2cm]
\answerExplain{$\frac{1}{9}$}{$\left(\frac{1}{3}\right)^2 = \frac{1}{3} \times \frac{1}{3} = \frac{1}{9}$.}
\end{multicols}
\end{practiceBox}

% ── Practice A: Product Rule ─────────────────────────────────────────────────
\begin{practiceBox}{Product Rule: $a^m \cdot a^n = a^{m+n}$}

\practiceHeader{Simplify. Write as a single power, then evaluate if possible.}

\begin{multicols}{2}
\prob $6^2 \cdot 6^3 =$ \answerBlank[3cm]
\answerExplain{$6^5 = 7{,}776$}{Add exponents (same base): $6^{2+3} = 6^5 = 7{,}776$.}

\prob $4^3 \cdot 4^2 =$ \answerBlank[3cm]
\answerExplain{$4^5 = 1{,}024$}{$4^{3+2} = 4^5 = 1{,}024$.}

\prob $7^1 \cdot 7^4 =$ \answerBlank[3cm]
\answerExplain{$7^5 = 16{,}807$}{$7^{1+4} = 7^5 = 16{,}807$.}

\prob $2^6 \cdot 2^2 =$ \answerBlank[3cm]
\answerExplain{$2^8 = 256$}{$2^{6+2} = 2^8 = 256$.}

\prob $\left(\frac{1}{2}\right)^2 \cdot \left(\frac{1}{2}\right)^3 =$ \answerBlank[3cm]
\answerExplain{$\left(\frac{1}{2}\right)^5 = \frac{1}{32}$}{Add exponents: $\left(\frac{1}{2}\right)^{2+3} = \left(\frac{1}{2}\right)^5 = \frac{1}{32}$.}
\end{multicols}
\end{practiceBox}

% ── Practice B: Quotient Rule ────────────────────────────────────────────────
\begin{practiceBox}[funTeal]{Quotient Rule: $\dfrac{a^m}{a^n} = a^{m-n}$}

\practiceHeader[funTeal]{Simplify. Write as a single power, then evaluate if possible.}

\begin{multicols}{2}
\prob $\dfrac{10^7}{10^4} =$ \answerBlank[3cm]
\answerExplain{$10^3 = 1{,}000$}{Subtract exponents: $10^{7-4} = 10^3 = 1{,}000$.}

\prob $\dfrac{5^6}{5^2} =$ \answerBlank[3cm]
\answerExplain{$5^4 = 625$}{$5^{6-2} = 5^4 = 625$.}

\prob $\dfrac{3^5}{3^3} =$ \answerBlank[3cm]
\answerExplain{$3^2 = 9$}{$3^{5-3} = 3^2 = 9$.}

\prob $\dfrac{8^4}{8^1} =$ \answerBlank[3cm]
\answerExplain{$8^3 = 512$}{$8^{4-1} = 8^3 = 512$.}

\prob $\dfrac{2^{10}}{2^7} =$ \answerBlank[3cm]
\answerExplain{$2^3 = 8$}{$2^{10-7} = 2^3 = 8$.}
\end{multicols}
\end{practiceBox}

% ── Practice C: Power of a Power ─────────────────────────────────────────────
\begin{practiceBox}[funOrange]{Power of a Power: $(a^m)^n = a^{mn}$}

\prob $(2^3)^4 =$ \answerBlank[3cm]
\answerExplain{$2^{12} = 4{,}096$}{Multiply exponents: $2^{3 \times 4} = 2^{12} = 4{,}096$.}

\prob $(5^2)^3 =$ \answerBlank[3cm]
\answerExplain{$5^6 = 15{,}625$}{$5^{2 \times 3} = 5^6 = 15{,}625$.}

\prob $(10^3)^2 =$ \answerBlank[3cm]
\answerExplain{$10^6 = 1{,}000{,}000$}{$10^{3 \times 2} = 10^6 = 1{,}000{,}000$.}

\prob $(3^1)^5 =$ \answerBlank[3cm]
\answerExplain{$3^5 = 243$}{$3^{1 \times 5} = 3^5 = 243$.}
\end{practiceBox}

% ── Practice D: Zero Exponent and Mixed ──────────────────────────────────────
\begin{practiceBox}[funPurple]{Zero Exponent and Mixed Practice}

\prob $(-3)^0 =$ \answerBlank[2cm]
\answerExplain{$1$}{Any nonzero number raised to the zero power equals $1$.}

\prob $100^0 =$ \answerBlank[2cm]
\answerExplain{$1$}{$100^0 = 1$ by the zero exponent rule.}

\prob $\dfrac{4^5 \cdot 4^2}{4^3} =$ \answerBlank[3cm]
\answerExplain{$4^4 = 256$}{Numerator: $4^{5+2} = 4^7$. Then $\frac{4^7}{4^3} = 4^{7-3} = 4^4 = 256$.}

\prob $(2^2)^3 \cdot 2^1 =$ \answerBlank[3cm]
\answerExplain{$2^7 = 128$}{First: $(2^2)^3 = 2^6$. Then $2^6 \cdot 2^1 = 2^{6+1} = 2^7 = 128$.}
\end{practiceBox}

% ── True or False ────────────────────────────────────────────────────────────
\begin{practiceBox}[funRed]{True or False?}

\trueOrFalse{$2^3 \cdot 3^2$ can be simplified to $6^5$ using the product rule.}
\answerExplain{False}{The product rule requires the same base. Since $2$ and $3$ are different bases, you cannot combine them. $2^3 \cdot 3^2 = 8 \times 9 = 72$.}

\trueOrFalse{$(4^3)^2 = 4^6$}
\answerExplain{True}{Power of a power: multiply exponents: $4^{3 \times 2} = 4^6$. This equals $4{,}096$.}
\end{practiceBox}

% ── Challenge ────────────────────────────────────────────────────────────────
\begin{challengeBox}
\prob Simplify $\dfrac{(3^2)^4 \cdot 3^3}{3^5}$ to a single power, then evaluate. \answerBlank[3cm]
\answerExplain{$3^6 = 729$}{$(3^2)^4 = 3^8$. Numerator: $3^8 \cdot 3^3 = 3^{11}$. Then $\frac{3^{11}}{3^5} = 3^{11-5} = 3^6 = 729$.}
\end{challengeBox}

\encouragement{Amazing work with exponent rules! These laws are the building blocks for algebra and beyond!}
